% Part: proof-theory
% Chapter: cut-elimination
% Section: introduction

\documentclass[../../../include/open-logic-section]{subfiles}

\begin{document}

\olfileid[id]{pt}{cut}{int}

\olsection{Pengantar}

Aturan~\Cut{} adalah aturan
\[
\Axiom$\Gamma \fCenter \Delta, !A$
\Axiom$!A, \Pi \fCenter \Lambda$
\RightLabel{\Cut}
\BinaryInf$\Gamma, \Pi \fCenter \Delta, \Lambda$
\DisplayProof
\]
(!!{formula}~$!A$ disebut \emph{formula potongan}.)

\Cut{} disertakan dalam kalkulus sekuen asli Gentzen,~\Log{LK}.
\emph{Hauptsatz} Gentzen (``teorema utama'') menyatakan bahwa setiap sekuen
yang !!{provable} dalam~\Log{LK} juga !!{provable} tanpa menggunakan
aturan~\Cut{}, yakni bahwa \Cut{} dapat ``dieliminasi'' dari bukti-bukti
\Log{LK}. Sistem kita, \Log{G1c} dan \Log{G3c}, tidak memuat \Cut, tetapi
pertanyaan yang sama dapat diajukan untuk keduanya: dapatkah \Cut{}
dieliminasi dari !!{proof}s yang menggunakan \Cut{} bersama aturan-aturan
lain \Log{G1c} (atau aturan-aturan \Log{G3c})? Dalam terminologi yang telah
kita tetapkan, pertanyaan yang ekuivalen ialah: Apakah \Cut{} merupakan
aturan yang dapat diterima dalam \Log{G1c} (atau \Log{G3c})?

Teorema eliminasi potongan barangkali merupakan hasil teori bukti yang paling
penting dan paling tersohor. Pertama-tama, teorema itu menetapkan bahwa
\Log{G1c} dan \Log{G3c} tidak memerlukan aturan yang sepintas mungkin tampak
jelas-jelas dibutuhkan. Ketika mencoba memformalkan bukti yang sebenarnya,
Anda akan mendapati bahwa diperlukan cara untuk menangani penggunaan lema.
Matematikawan suka membuktikan hasil, kemudian mengacu pada hasil tersebut
dalam pembuktian hasil lain. Jadi, mula-mula Anda dapat memformalkan bukti
hasil~$L$ (lema) sebagai !!{proof} sekuen $\Sequent L$. Lalu Anda
memformalkan bukti hasil kedua~$R$, yang menggunakan~$L$, sebagai !!{proof}
sekuen $L \Sequent R$. Bagaimana Anda tahu bahwa terdapat bukti~$R$, yakni
!!a{proof} dari $\Sequent R$ saja? Anda memerlukan cara untuk menggabungkan
kedua !!{proof} itu, dan aturan~\Cut{} memungkinkan Anda melakukannya.

Kita telah menyebutkan bahwa \Log{G1c} dan \Log{G3c} lengkap, sehingga
keduanya membuktikan setiap sekuen yang sewajarnya diharapkan dapat
dibuktikan---semua sekuen yang valid. Hal ini dapat dibuktikan secara
langsung. Namun, ketika Gentzen menulis, hanya sistem derivasi aksiomatik
yang diketahui lengkap. Untuk menunjukkan bahwa \Log{LK} lengkap, Gentzen
memberikan terjemahan !!{derivation}s dalam sistem aksiomatik menjadi
!!{proof}s \Log{LK}. Terjemahan ini menggunakan \Cut{} secara esensial.
Eliminasi potongan bukan hanya hasil penting bagi logika klasik: banyak
logika lain memiliki kalkulus sekuen. Bagi beberapa logika, bukti langsung
kelengkapan kalkulus sekuennya sulit atau mustahil, sehingga terjemahan dari
sistem lain (yang menggunakan \Cut) merupakan satu-satunya cara untuk
menunjukkan kelengkapan. Karena itu, pembuktian eliminasi potongan secara
teori-bukti diperlukan untuk menetapkan bahwa \Cut{} berlebih dan bahwa
sistem tanpa \Cut{} lengkap.

Eliminasi potongan juga penting karena dapat diterapkan untuk memperoleh
hasil-hasil lain yang menarik secara mandiri. Dua kasus yang akan kita
bahas ialah lema interpolasi Craig dan teorema Herbrand.

Gagasan yang digunakan untuk membuktikan eliminasi potongan biasanya
mencakup cara menunjukkan bagaimana !!a{proof} yang menggunakan \Cut{}
dapat ditransformasikan menjadi bukti tanpa aturan tersebut. Transformasi
ini lazimnya ``lokal'', yakni kita mengambil suatu bagian !!a{proof}
(biasanya sub-!!{proof} yang berakhir pada inferensi \Cut{}) dan
menggantinya dengan sub-!!{proof} lain dari sekuen yang sama, tetapi
inferensi \Cut{} di dalamnya telah dihapus atau diganti oleh inferensi lain.
Argumen semacam ini menyediakan algoritme langkah demi langkah untuk
mentransformasikan !!{proof}s dengan \Cut{} menjadi bukti tanpa potongan.
Hal ini memberikan lebih banyak informasi ketika prosedur eliminasi
potongan digunakan dalam penerapan. Sebagai contoh, terdapat bukti
teori-model atas lema interpolasi yang berbentuk ``jika $!A \lif !B$ valid,
maka suatu interpolan ada'' (untuk sementara, abaikan apa itu interpolan).
Argumen teori-bukti yang menggunakan eliminasi potongan menyediakan
algoritme yang menerima !!a{proof} dari $!A \lif !B$ dan mengembalikan
sebuah interpolan. Berbeda dengan argumen teori-model, argumen teori-bukti
tersebut bersifat \emph{konstruktif}.

Ada dua cara populer untuk membuktikan teorema eliminasi potongan. Cara
pertama (yang berasal dari Gentzen) menunjukkan bahwa kita dapat menghapus
potongan-potongan teratas dalam suatu bukti, lalu terus menghapusnya sampai
tidak ada yang tersisa. Cara kedua (mula-mula dikemukakan oleh Sch\"utte dan
Tait) berfokus pada potongan-potongan paling kompleks dalam !!a{proof}, dan
secara berturut-turut menghapus potongan yang maksimal dalam arti tidak ada
potongan lain dengan !!{formula} potongan berkedalaman lebih tinggi. Masing-masing
memiliki kelebihan dan kekurangan. Bagi beberapa sistem, satu pendekatan
berhasil sedangkan pendekatan lain sulit diterapkan atau bahkan mustahil.
Karena itu, kita menguraikan kedua pendekatan tersebut. Kita akan membuktikan
eliminasi potongan untuk~\Log{G3c}. Bahkan, yang akan kita tunjukkan bukanlah
bahwa \Cut{} dapat dieliminasi, melainkan \emph{potongan yang berbagi
konteks}~\CutCS:
\[
\Axiom$\Gamma \fCenter \Delta, !A$
\Axiom$!A, \Gamma \fCenter \Delta$
\RightLabel{\Cut}
\BinaryInf$\Gamma \fCenter \Delta$
\DisplayProof
\]
Jika suatu kalkulus sekuen mengizinkan pelemahan dan kontraksi (baik sebagai
aturan nyata seperti dalam \Log{G1c} maupun sebagai aturan yang dapat
diterima seperti dalam \Log{G3c}), \Cut{} dapat menyimulasikan \CutCS, dan
sebaliknya. Kita memilih \CutCS{} alih-alih \Cut karena strategi eliminasi
potongan pertama lebih mudah diuraikan untuk~\CutCS, sedangkan strategi
kedua hanya berlaku untuk~\CutCS.

\end{document}
